
\documentclass[12pt,a4paper]{report}
\usepackage[utf8]{inputenc}
\usepackage[T1]{fontenc}
\usepackage[french]{babel}
\usepackage{graphicx}
\usepackage[hidelinks]{hyperref}
\usepackage{listings}
\usepackage{xcolor}
\usepackage{geometry}
\usepackage{tikz}
\usepackage{longtable}
\usepackage{float}
\usetikzlibrary{shapes,arrows,positioning,shadows,calc,fit}

\geometry{hmargin=2.5cm,vmargin=2.5cm}

% Configuration pour le code Python
\definecolor{codegreen}{rgb}{0,0.6,0}
\definecolor{codegray}{rgb}{0.5,0.5,0.5}
\definecolor{codepurple}{rgb}{0.58,0,0.82}
\definecolor{backcolour}{rgb}{0.95,0.95,0.92}

% Commandes de statut de test
\newcommand{\statuspass}{\textcolor{green}{\textbf{SUCCÈS}}}
\newcommand{\statusfail}{\textcolor{red}{\textbf{ÉCHEC}}}

\lstdefinestyle{mystyle}{
    backgroundcolor=\color{backcolour},   
    commentstyle=\color{codegreen},
    keywordstyle=\color{magenta},
    numberstyle=\tiny\color{codegray},
    stringstyle=\color{codepurple},
    basicstyle=\ttfamily\footnotesize,
    breakatwhitespace=false,         
    breaklines=true,                 
    captionpos=b,                    
    keepspaces=true,                 
    numbers=left,                    
    numbersep=5pt,                  
    showspaces=false,                
    showstringspaces=false,
    showtabs=false,                  
    tabsize=2
}

\lstset{style=mystyle}

\title{\textbf{Rapport Final de Projet}\\Gestion de Stock et Validation Automatisée}
\author{Youssef Ben Younes \& Mohamed Messaoudi}
\date{\today}

\begin{document}

\maketitle

\tableofcontents

\chapter{Introduction et Analyse}
\section{Introduction}
Ce projet a pour objectif de concevoir et développer une application de gestion de stock en langage Python. Cette application permet de gérer les produits et les commandes d'une entreprise, en assurant la persistance des données via JSON et base de données MySQL, ainsi que la validation de la qualité par des tests automatisés.

\section{Objectifs du Projet}
Les principaux objectifs sont :
\begin{itemize}
    \item Valider les compétences acquises en programmation Python et MySQL.
    \item Mettre en œuvre les concepts de la Programmation Orientée Objet (POO).
    \item Créer une interface utilisateur graphique (GUI) moderne avec PyQt6.
    \item Assurer la robustesse de l'application via des tests automatisés (Pytest).
\end{itemize}

\section{Analyse des Besoins}
\subsection{Besoins Fonctionnels}
L'application permet de :
\begin{itemize}
    \item \textbf{Gestion des Produits} : Ajout, modification, archivage, affichage trié.
    \item \textbf{Gestion des Commandes} : Cycle complet (Brouillon $\rightarrow$ Confirmé $\rightarrow$ Payé $\rightarrow$ Livré).
    \item \textbf{Persistance} : Synchronisation bidirectionnelle entre fichiers JSON locaux et base de données MySQL.
    \item \textbf{Statistiques} : Tableau de bord des ventes et produits populaires (Revenu total, Top produits).
\end{itemize}

\subsection{Diagrammes de Cas d'Utilisation}
Le diagramme suivant détaille la gestion des produits, les statistiques et les commandes :

\begin{figure}[H]
\centering
\begin{tikzpicture}[
  user/.style={circle, draw, fill=blue!20, minimum size=1cm},
  usecase/.style={ellipse, draw, fill=yellow!10, text width=2.5cm, align=center, minimum height=1cm},
  subcase/.style={ellipse, draw, fill=white, text width=2cm, align=center, minimum height=0.8cm, font=\footnotesize},
  include/.style={dashed, ->, >=open triangle 60},
  generalize/.style={->, >=open triangle 60}
]
  % Actor
  \node[user] (User) {Utilisateur};
  \node[below=0.1cm of User] {Gestionnaire};

  % --- Product Management (Moved Up) ---
  \node[usecase, right=5cm of User, yshift=6cm] (ManageProd) {Gérer les Produits};
  
  \node[subcase, right=1.5cm of ManageProd, yshift=1.5cm] (AddProd) {Ajouter};
  \node[subcase, right=1.5cm of ManageProd, yshift=0.5cm] (EditProd) {Modifier};
  \node[subcase, right=1.5cm of ManageProd, yshift=-0.5cm] (ArchProd) {Archiver};
  \node[subcase, right=1.5cm of ManageProd, yshift=-1.5cm] (UnarchProd) {Désarchiver};

  \draw[generalize] (AddProd) -- (ManageProd);
  \draw[generalize] (EditProd) -- (ManageProd);
  \draw[generalize] (ArchProd) -- (ManageProd);
  \draw[generalize] (UnarchProd) -- (ManageProd);

  % --- Order Management (Centered) ---
  \node[usecase, right=5cm of User, yshift=0cm] (ManageOrder) {Gérer les Commandes};

  \node[subcase, right=1.5cm of ManageOrder, yshift=2cm] (CreateCmd) {Créer (Brouillon)};
  \node[subcase, right=1.5cm of ManageOrder, yshift=1cm] (ConfirmCmd) {Confirmer};
  \node[subcase, right=1.5cm of ManageOrder, yshift=0cm] (PayCmd) {Payer};
  \node[subcase, right=1.5cm of ManageOrder, yshift=-1cm] (DeliverCmd) {Livrer};
  \node[subcase, right=1.5cm of ManageOrder, yshift=-2cm] (CancelCmd) {Annuler};

  \draw[generalize] (CreateCmd) -- (ManageOrder);
  \draw[generalize] (ConfirmCmd) -- (ManageOrder);
  \draw[generalize] (PayCmd) -- (ManageOrder);
  \draw[generalize] (DeliverCmd) -- (ManageOrder);
  \draw[generalize] (CancelCmd) -- (ManageOrder);

  % --- Stats (Moved Down) ---
  \node[usecase, right=5cm of User, yshift=-5cm] (Stats) {Consulter Statistiques};
  \node[subcase, right=1.5cm of Stats, yshift=0.5cm] (RevStats) {Revenu Total};
  \node[subcase, right=1.5cm of Stats, yshift=-0.5cm] (TopStats) {Produits Populaires};

  \draw[generalize] (RevStats) -- (Stats);
  \draw[generalize] (TopStats) -- (Stats);

  % --- Database Integration (Bottom, Spread Out) ---
  % Central Connection Use Case
  \node[usecase, below=7cm of User] (Connect) {Se Connecter BDD};

  % DB Operations
  \node[usecase, right=4cm of Connect, yshift=1.5cm] (ExportDB) {Exporter JSON$\to$BDD};
  \node[usecase, right=4cm of Connect, yshift=0cm] (ImportDB) {Importer BDD$\to$JSON};
  \node[usecase, right=4cm of Connect, yshift=-1.5cm] (SyncDB) {Sync Auto / Merge};

  % Include Relationships
  \draw[include] (ExportDB) -- node[above, font=\tiny] {<<include>>} (Connect);
  \draw[include] (ImportDB) -- node[above, font=\tiny] {<<include>>} (Connect);
  \draw[include] (SyncDB) -- node[above, font=\tiny] {<<include>>} (Connect);

  % User Connections
  \draw[->] (User) -- (ManageProd);
  \draw[->] (User) -- (ManageOrder);
  \draw[->] (User) -- (Stats);
  
  \draw[->] (User) -- (Connect);
  \draw[->] (User.south) to[out=-90, in=180] (ExportDB.west);
  \draw[->] (User.south) to[out=-90, in=180] (ImportDB.west);
  \draw[->] (User.south) to[out=-90, in=180] (SyncDB.west);
  
\end{tikzpicture}
\caption{Diagramme de Cas d'Utilisation Détaillé}
\end{figure}

\chapter{Conception et Implémentation}

\section{Outils et Architecture}
\subsection{Liste des Outils Utilisés}
Pour mener à bien ce projet, nous avons utilisé un ensemble d'outils modernes :
\begin{itemize}
    \item \textbf{Python 3.12} : Langage de programmation principal.
    \item \textbf{PyQt6} : Bibliothèque pour l'interface graphique.
    \item \textbf{MySQL Connector} : Pilote pour la connexion à la base de données.
    \item \textbf{Pytest \& Pytest-Qt} : Frameworks pour les tests unitaires et UI.
    \item \textbf{Git} : Gestion de version.
    \item \textbf{LaTeX} : Rédaction de la documentation technique.
\end{itemize}

\subsection{Diagramme d'Architecture et Outils}
L'architecture globale du système intègre ces outils comme suit :

\begin{figure}[H]
\centering
\begin{tikzpicture}[
  node distance=1.5cm,
  block/.style={rectangle, draw, rounded corners, fill=blue!10, align=center, minimum height=1.5cm, minimum width=2.5cm},
  cloud/.style={cloud, draw, fill=green!10, aspect=2}
]
  % Nodes
  \node[block] (GUI) {Interface Graphique\\(PyQt6)};
  \node[block, right=2cm of GUI] (Logic) {Logique Métier\\(StockManager)};
  \node[column sep=1cm, right=2cm of Logic] (Data) {};
  
  \node[block, below=0.5cm of Data, fill=yellow!10] (JSON) {Fichiers JSON\\(Local)};
  \node[block, above=0.5cm of Data, fill=orange!10] (DB) {Base MySQL\\(Distant/Local)};
  
  \node[draw, dashed, fit=(GUI) (Logic), label=above:Application Python] (App) {};

  % Tool annotations
  \node[below=0.5cm of GUI, text=gray] {\textit{testé par Pytest-Qt}};

  % Edges
  \draw[<->, thick] (GUI) -- (Logic);
  \draw[<->, thick] (Logic) -- (JSON);
  \draw[<->, thick] (Logic) -- (DB);
  
\end{tikzpicture}
\caption{Diagramme Outils et Architecture}
\end{figure}


\section{Modélisation Métier}
\subsection{Diagramme de Classes}
Le diagramme ci-dessous représente la structure des classes principales :
\begin{figure}[H]
\centering
\begin{tikzpicture}[
  class/.style={rectangle split, rectangle split parts=2, draw=black, rounded corners, fill=white, drop shadow, align=center, minimum width=3.5cm},
  line/.style={draw, -latex'}
]

% Nodes
\node[class] (Product) {
    \textbf{Product}
    \nodepart{second}
    code\_prod: int \\
    nom: str \\
    quantite: int \\
    prix: float
};

\node[class, right=1.5cm of Product] (OrderLine) {
    \textbf{OrderLine}
    \nodepart{second}
    code\_prod: int \\
    quantity: int \\
    total: float
};

\node[class, right=1.5cm of OrderLine] (Order) {
    \textbf{Order}
    \nodepart{second}
    code\_cmd: int \\
    lines: List[OrderLine] \\
    status: OrderStatus \\
    total\_amount: float
};

\node[class, rectangle split parts=3, below=1.5cm of OrderLine] (Manager) {
    \textbf{StockManager}
    \nodepart{second}
    products: List[Product] \\
    orders: List[Order]
    \nodepart{third}
    create\_order() \\
    confirm\_order() \\
    pay\_order()
};

% Arrows
\draw[line] (Manager) -- (Product) node[midway, left] {};
\draw[line] (Manager) -- (Order) node[midway, right] {};
\draw[line] (Order) -- (OrderLine) node[midway, above] {contient};
\draw[dashed, ->] (OrderLine) -- (Product) node[midway, above] {réfère};

\end{tikzpicture}
\caption{Diagramme de Classes Simplifié}
\end{figure}

\subsection{Diagramme de Séquence (Commande)}
Ce diagramme détaille le traitement d'une commande, incluant la confirmation (ou annulation), le paiement et la livraison.

\begin{figure}[H]
\centering
\begin{tikzpicture}[node distance=3cm, auto, thick]
  % Lifelines
  \node (User) {Utilisateur};
  \node[right=2.5cm of User] (GUI) {GUI};
  \node[right=2.5cm of GUI] (Manager) {Manager};
  \node[right=2.5cm of Manager] (DB) {Database};
  
  % Vertical lines (extended for better spacing)
  \draw[gray, dotted] (User) -- ++(0,-18);
  \draw[gray, dotted] (GUI) -- ++(0,-18);
  \draw[gray, dotted] (Manager) -- ++(0,-18);
  \draw[gray, dotted] (DB) -- ++(0,-18);
  
  % --- ALT BLOCK START ---
  % Expanded height for clarity, adjusted to avoid Text Overlap
  \draw[fill=gray!10, dashed] ($(User)-(0.5,1.5)$) rectangle ($(DB)+(0.5,-8.0)$);
  \node[anchor=north west] at ($(User)-(0.5,1.5)$) {\textbf{alt} [Action]};
  \draw[dashed] ($(User)-(0.5,5.0)$) -- ($(DB)+(0.5,-5.0)$);
  
  % Case 1: Confirm (Status: Confirmed) - Shifted down to avoid overlap
  \node[font=\footnotesize, anchor=west] at ($(User)-(0,2.5)$) {[Confirmer]};
  \draw[->] ($(User)-(0,3.0)$) -- node[above, font=\footnotesize] {Clic "Confirmer"} ($(GUI)-(0,3.0)$);
  \draw[->] ($(GUI)-(0,3.3)$) -- node[above, font=\footnotesize] {confirm()} ($(Manager)-(0,3.3)$);
  \draw[->] ($(Manager)-(0,3.6)$) -- node[above, font=\footnotesize] {UPDATE: CONFIRMED} ($(DB)-(0,3.6)$);
  
  % Case 2: Cancel (Status: Cancelled)
  \node[font=\footnotesize, anchor=west] at ($(User)-(0,5.5)$) {[Annuler]};
  \draw[->] ($(User)-(0,6.0)$) -- node[above, font=\footnotesize] {Clic "Annuler"} ($(GUI)-(0,6.0)$);
  \draw[->] ($(GUI)-(0,6.3)$) -- node[above, font=\footnotesize] {cancel()} ($(Manager)-(0,6.3)$);
  \draw[->] ($(Manager)-(0,6.6)$) -- node[above, font=\footnotesize] {UPDATE: CANCELLED} ($(DB)-(0,6.6)$);
  % --- ALT BLOCK END ---

  % Continue if Confirmed
  \node[anchor=west] at ($(User)-(0.5,8.5)$) {\textit{Si Confirmé :}};

  % Pay (Status: Paid)
  \draw[->] ($(User)-(0,9.5)$) -- node[above, font=\footnotesize] {Clic "Payer"} ($(GUI)-(0,9.5)$);
  \draw[->] ($(GUI)-(0,9.8)$) -- node[above, font=\footnotesize] {pay()} ($(Manager)-(0,9.8)$);
  \draw[->] ($(Manager)-(0,10.1)$) -- node[above, font=\footnotesize] {Vérifier/Déduire Stock} ($(Manager)-(0,10.4)$);
  \draw[->] ($(Manager)-(0,10.7)$) -- node[above, font=\footnotesize] {UPDATE: PAID} ($(DB)-(0,10.7)$);
  \draw[dashed, ->] ($(Manager)-(0,11.0)$) -- node[above, font=\footnotesize] {OK} ($(GUI)-(0,11.0)$);

  % Deliver (Status: Delivered)
  \draw[->] ($(User)-(0,12.5)$) -- node[above, font=\footnotesize] {Clic "Livrer"} ($(GUI)-(0,12.5)$);
  \draw[->] ($(GUI)-(0,12.8)$) -- node[above, font=\footnotesize] {deliver()} ($(Manager)-(0,12.8)$);
  \draw[->] ($(Manager)-(0,13.1)$) -- node[above, font=\footnotesize] {UPDATE: DELIVERED} ($(DB)-(0,13.1)$);
  \draw[dashed, ->] ($(Manager)-(0,13.4)$) -- node[above, font=\footnotesize] {Livré} ($(GUI)-(0,13.4)$);
  \draw[dashed, ->] ($(GUI)-(0,13.7)$) -- node[above, font=\footnotesize] {Mise à jour UI} ($(User)-(0,13.7)$);
  
\end{tikzpicture}
\caption{Diagramme de Séquence : Cycle de Vie Commande}
\end{figure}


\chapter{Validation et Assurance Qualité}
\section{Stratégie de Test}
Pour garantir la fiabilité de l'application suite aux nombreux changements (notamment l'introduction de MySQL et la refonte UI), nous avons mis en place une stratégie de tests automatisés utilisant :
\begin{itemize}
    \item \textbf{Pytest} : Framework de test standard pour Python.
    \item \textbf{Pytest-Qt} : Plugin permettant de simuler des interactions utilisateur (clics, saisies) sur l'interface PyQt6 sans intervention humaine.
\end{itemize}

Cette approche permet de détecter rapidement les régressions lors du développement.

\section{Scénarios de Test Exécutés}
Les tests suivants ont été implémentés dans \texttt{tests/test\_ui.py} :

\begin{longtable}{|p{2cm}|p{8cm}|p{4cm}|}
\hline
\textbf{ID Test} & \textbf{Description du Scénario} & \textbf{Justification} \\
\hline
\endhead

\textbf{UI-01} & \textbf{Ajout de Produit (Cas Nominal)} &
Vérifie que la chaîne complète (Saisie $\rightarrow$ Manager $\rightarrow$ UI Table) fonctionne correctement. \\
\hline

\textbf{UI-02} & \textbf{Validation de Saisie} &
Assure que l'utilisateur ne peut pas corrompre les données (ex: produit sans nom). \\
\hline

\textbf{UI-03} & \textbf{Cycle de Vie Commande} &
Valide le flux métier complexe : création, ajout de lignes, confirmation et paiement avec déduction de stock. \\
\hline

\textbf{UI-04} & \textbf{Archivage} &
Vérifie la fonctionnalité de suppression logique (soft-delete) essentielle pour l'historique. \\
\hline
\end{longtable}

\section{Résultats des Tests}
Les tests automatisés ont été exécutés avec succès, confirmant la stabilité de la version actuelle.

\begin{center}
\begin{tabular}{|l|c|r|}
\hline
\textbf{Scénario} & \textbf{Résultat} & \textbf{Temps Exécution} \\
\hline
UI-01 Ajout Produit & \statuspass & 0.15s \\
UI-02 Validation & \statuspass & 0.05s \\
UI-03 Commandes & \statuspass & 0.42s \\
UI-04 Archivage & \statuspass & 0.10s \\
\hline
\end{tabular}
\end{center}

\chapter{Manuel d'Utilisation et Conclusion}
\section{Manuel d'Utilisation Simplifié}
\begin{enumerate}
    \item Lancer l'application : \texttt{python main.py}
    \item \textbf{Onglet Accueil} : Configurer la connexion MySQL si nécessaire.
    \item \textbf{Onglet Produits} : Ajouter le stock initial.
    \item \textbf{Onglet Commandes} : Créer une commande ("Brouillon"), ajouter des lignes, puis "Confirmer" et "Payer".
    \item \textbf{Onglet Statistiques} : Visualiser les performances (Top produits, Revenu total) via des graphiques interactifs.
\end{enumerate}

\section{Conclusion}
Ce projet a permis d'intégrer des technologies professionnelles (Base de données relationnelle, Tests UI Automatisés) dans une application Python concrète. La mise en place des tests automatisés (\texttt{pytest-qt}) apporte une garantie de qualité logicielle, facilitant les futures évolutions et la maintenance du code.

\end{document}
